\chapter{METODOLOGI PENELITIAN}

\petunjuk{Jelaskan objek penelitian, data, alat, tahapan, metode pengumpulan data, metode analisis, serta cara pengujian. Bab ini juga memuat contoh persamaan dan kode program.}

\section{Tahapan Penelitian}

Tahapan penelitian harus disajikan secara berurutan dan dapat direplikasi. Contoh tahapan umum adalah identifikasi masalah, pengumpulan data, perancangan solusi, implementasi, pengujian, dan evaluasi.

\section{Contoh Persamaan}

Sebelum persamaan ditampilkan, jelaskan variabel dan tujuan penggunaannya. Nilai rata-rata dapat dihitung menggunakan Persamaan~\eqref{eq:rata-rata}.

\begin{equation}
  \bar{x} = \frac{1}{n}\sum_{i=1}^{n}x_i
  \label{eq:rata-rata}
\end{equation}

Pada Persamaan~\eqref{eq:rata-rata}, $\bar{x}$ adalah nilai rata-rata, $x_i$ adalah nilai pengamatan ke-$i$, dan $n$ adalah jumlah pengamatan. Nomor persamaan dibuat otomatis berdasarkan nomor bab.

\section{Contoh Kode Program}

Notasi algoritmik, kode program, atau pseudocode menggunakan font monospaced 9 pt dan spasi tunggal. Kode yang panjang sebaiknya ditempatkan di lampiran, sedangkan bagian inti hanya menampilkan potongan yang dibahas. Contoh ditunjukkan pada Gambar~\ref{fig:contoh-kode}.

\begin{figure}[H]
\begin{minipage}{\textwidth}
\begin{lstlisting}[language=Python]
def hitung_rata_rata(data):
    if not data:
        return 0
    return sum(data) / len(data)
\end{lstlisting}
\end{minipage}
\caption{Contoh kode program untuk menghitung nilai rata-rata}
\label{fig:contoh-kode}
\end{figure}

\section{Metode Pengujian}

Jelaskan skenario, data uji, metrik, kriteria keberhasilan, dan cara interpretasi hasil. Setiap klaim pada bab hasil harus dapat ditelusuri kembali ke metode yang dijelaskan pada bagian ini.
